\chapter{Mixed Hodge Structures}\label{mixed-hodge-structures}
\chapterstatus{complete}

\section{Introduction}\label{mixed-hodge-structures:section-introduction}

The cohomology of a smooth but non-compact variety, or of a singular projective one, is not a pure Hodge structure. For example, $H^1(\CC^*; \QQ)$ is one-dimensional, so it cannot carry a
Hodge structure of odd weight. Deligne showed that the cohomology of every complex algebraic variety carries a \emph{mixed Hodge structure}: an increasing \emph{weight filtration} whose graded pieces
are pure Hodge structures of various weights. Morphisms of mixed Hodge structures are strict for both filtrations (Theorem~\ref{mixed-hodge-structures:theorem-abelian}), and this rigidity has consequences
for compact varieties too. Three of them are proved in Section~\ref{mixed-hodge-structures:section-applications}: the global invariant cycle theorem, used in Chapter~\ref{vhs}; the fact that
algebraicity of a Hodge class can be tested on a Zariski open subset, modulo lower dimensions; and the birational invariance of the Hodge conjecture, modulo lower dimensions. References:
\cite{deligne-hodge-2}, \cite{deligne-hodge-3}, \cite{peters-steenbrink}, \cite[Chapter 8]{voisin-hodge-1}.

\noindent\textbf{Prerequisites.} Chapter~\ref{hodge-structures} (Hodge Structures), Chapter~\ref{derived-categories} (Derived Categories) and Chapter~\ref{gaga} (Analytic Spaces and GAGA).

\noindent\textbf{Conventions.} Filtrations are finite and exhaustive. $W_\bullet$ always denotes an increasing filtration ($W_{k-1} \subseteq W_k$) and $F^\bullet$ a decreasing one
($F^{p+1} \subseteq F^p$). For an increasing filtration, $\mathrm{Gr}^W_k = W_k/W_{k-1}$. The Tate structure $\QQ(m)$ is $(2\pi i)^m\QQ \subseteq \CC$, pure of weight $-2m$ and type $(-m,-m)$
(Example~\ref{hodge-structures:example-hodge-structures}), and $H(m) = H \otimes \QQ(m)$. Varieties are complex algebraic varieties, and their cohomology is the singular cohomology of the associated
analytic space.

\section{Mixed Hodge structures}\label{mixed-hodge-structures:section-definition}

\begin{definition}\label{mixed-hodge-structures:definition-mhs}
A \emph{mixed Hodge structure} $(H, W, F)$ consists of
\begin{enumerate}
\item a finite-dimensional $\QQ$-vector space $H$;
\item a finite increasing filtration $W_\bullet$ of $H$ by $\QQ$-subspaces, the \emph{weight filtration};
\item a finite decreasing filtration $F^\bullet$ of $H_\CC = H \otimes_\QQ \CC$ by $\CC$-subspaces, the \emph{Hodge filtration};
\end{enumerate}
such that for every $k$ the filtration induced by $F$ on $\mathrm{Gr}^W_kH_\CC = W_kH_\CC/W_{k-1}H_\CC$, namely
\[
F^p\mathrm{Gr}^W_kH_\CC = (F^p \cap W_kH_\CC + W_{k-1}H_\CC)/W_{k-1}H_\CC,
\]
is the Hodge filtration of a pure $\QQ$-Hodge structure of weight $k$ on $\mathrm{Gr}^W_kH$ (Proposition~\ref{hodge-structures:proposition-filtration}). A \emph{morphism}
$f \colon (H, W, F) \to (H', W', F')$ is a $\QQ$-linear map with $f(W_k) \subseteq W'_k$ and $f_\CC(F^p) \subseteq F'^p$ for all $k$, $p$. The \emph{Hodge numbers} are
$h^{p,q}(H) = \dim_\CC(\mathrm{Gr}^W_{p+q}H_\CC)^{p,q}$, and the \emph{weights} of $H$ are the $k$ with $\mathrm{Gr}^W_kH \neq 0$.
\end{definition}

\begin{example}\label{mixed-hodge-structures:example-mhs}
\begin{enumerate}
\item A pure Hodge structure of weight $k$ is a mixed Hodge structure with $W_{k-1} = 0$ and $W_k = H$, and morphisms of pure Hodge structures of the same weight are the same as morphisms of the
associated mixed Hodge structures. A direct sum of pure Hodge structures of different weights is a mixed Hodge structure, with $W_k$ the sum of the summands of weight $\leq k$.
\item $H^1(\CC^*; \QQ)$ is one-dimensional, spanned by the class of $\frac{1}{2\pi i}\frac{dz}{z}$, which integrates to $1$ over the unit circle. With $W_1 = 0$, $W_2 = H^1$, $F^1 = H^1(\CC^*; \CC)$ and
$F^2 = 0$, it is the pure Tate structure $\QQ(-1)$ of weight $2$ and type $(1,1)$. This is the structure given by Theorem~\ref{mixed-hodge-structures:theorem-smooth} below, with $\bar X = \PP^1$ and $D = \{0, \infty\}$.
\item Let $C$ be a nodal plane cubic curve, obtained from $\PP^1$ by identifying two points. Then $H^1(C; \QQ) \cong \QQ$ is pure of weight $0$ and type $(0,0)$
(Example~\ref{mixed-hodge-structures:example-nodal} below): the loop through the node is not detected by holomorphic $1$-forms.
\item \emph{Kummer structures.} Let $q \in \CC^*$, choose $\log q$, and put $\lambda = \log q/2\pi i$. On $H = \QQ e_0 \oplus \QQ e_1$ let $W_{-3} = 0$, $W_{-2} = W_{-1} = \QQ e_1$, $W_0 = H$, and
$F^{-1} = H_\CC$, $F^0 = \CC(e_0 - \lambda e_1)$, $F^1 = 0$. This is a mixed Hodge structure $K_q$, an extension of $\QQ(0)$ by $\QQ(1)$ (Exercise~\ref{mixed-hodge-structures:exercise-kummer}). Another choice
of $\log q$ changes $\lambda$ by an integer $n$, and $e_0 \mapsto e_0 + ne_1$ is an isomorphism between the two structures, so $K_q$ depends only on $q$.
\end{enumerate}
\end{example}

\begin{counterexample}\label{mixed-hodge-structures:counterexample-no-splitting}
A mixed Hodge structure need not be the direct sum of its graded pieces, not even after extending scalars to $\RR$. In $K_q$, a complement of $W_{-2}$ that is a sub-mixed Hodge structure is a line
$L = \QQ(e_0 + te_1)$ with $t \in \QQ$ whose induced Hodge filtration is that of $\QQ(0)$, that is, with $L_\CC \subseteq F^0 = \CC(e_0 - \lambda e_1)$. This forces $t = -\lambda$. So $K_q$ splits if and
only if $\lambda \in \QQ$, that is, $q$ is a root of unity; and the same argument with $t \in \RR$ shows that $K_q \otimes \RR$ does not split when $|q| \neq 1$, since then $\lambda \notin \RR$. Carlson showed
that the analogous extensions of $\ZZ(0)$ by $\ZZ(1)$ over $\ZZ$ are classified exactly by $q \in \CC/2\pi i\ZZ \cong \CC^*$ \cite{carlson-1980}: Kummer structures are the Hodge-theoretic incarnation of
the exponential map. Contrast this with pure polarisable Hodge structures, which form a semisimple category (Theorem~\ref{hodge-structures:theorem-semisimple}).
\end{counterexample}

\section{The Deligne splitting and strictness}\label{mixed-hodge-structures:section-strictness}

\begin{theorem}[Deligne]\label{mixed-hodge-structures:theorem-splitting}
Let $(H, W, F)$ be a mixed Hodge structure. For $p, q \in \ZZ$ put
\[
I^{p,q} = F^p \cap W_{p+q,\CC} \cap \Big(\overline{F^q} \cap W_{p+q,\CC} + \sum_{j \geq 1}\overline{F^{q-j}} \cap W_{p+q-j-1,\CC}\Big).
\]
Then $H_\CC = \bigoplus_{p,q}I^{p,q}$, and
\[
W_{k,\CC} = \bigoplus_{p+q \leq k}I^{p,q}, \qquad F^p = \bigoplus_{p' \geq p}I^{p',q}, \qquad \overline{I^{p,q}} \subseteq I^{q,p} \oplus \bigoplus_{p' < q,\,q' < p}I^{p',q'}.
\]
The projection $W_{p+q,\CC} \to \mathrm{Gr}^W_{p+q}H_\CC$ maps $I^{p,q}$ isomorphically onto $(\mathrm{Gr}^W_{p+q})^{p,q}$. Conversely, if $W$ is a filtration by $\QQ$-subspaces and $F$ a filtration of $H_\CC$
admitting a bigrading $H_\CC = \bigoplus I^{p,q}$ with these three properties, then $(H, W, F)$ is a mixed Hodge structure.
\end{theorem}

The first part is \cite{deligne-hodge-2}; see \cite[Chapter 3]{peters-steenbrink} for the explicit formula. The converse is immediate: on $\mathrm{Gr}^W_k$ the images of the $I^{p,q}$ with $p + q = k$
give a decomposition with $\overline{H^{p,q}} = H^{q,p}$, since the error terms have $p' + q' < k$, and the filtration induced by $F$ is $\bigoplus_{p' \geq p}H^{p',k-p'}$. The subspaces $I^{p,q}$ form
the \emph{Deligne splitting}. It splits $W$ and $F$ simultaneously, but it is not compatible with complex conjugation in general (Counterexample~\ref{mixed-hodge-structures:counterexample-no-splitting}).
Every morphism $f$ of mixed Hodge structures maps $I^{p,q}(H)$ into $I^{p,q}(H')$, as is clear from the formula, since $f$ preserves $W$, $F$ and, being defined over $\QQ$, also $\overline F$.

\begin{theorem}[Deligne]\label{mixed-hodge-structures:theorem-abelian}
\begin{enumerate}
\item Every morphism $f \colon H \to H'$ of mixed Hodge structures is \emph{strict} for $W$ and for $F$: $f(W_k) = W'_k \cap f(H)$ and $f_\CC(F^p) = F'^p \cap f_\CC(H_\CC)$.
\item Mixed Hodge structures form an abelian category, in which kernels, images and cokernels carry the induced filtrations.
\item For every $k$, the functor $H \mapsto \mathrm{Gr}^W_kH$ from mixed Hodge structures to pure Hodge structures of weight $k$ is exact.
\item $H \otimes H'$, $H^\vee$ and $\Hom(H, H') = H^\vee \otimes H'$ carry mixed Hodge structures, with
$W_k(H \otimes H') = \sum_{a+b=k}W_aH \otimes W_bH'$, $F^p(H \otimes H') = \sum_{a+b=p}F^aH \otimes F^bH'$, $W_k(H^\vee) = (W_{-k-1}H)^\perp$ and $F^p(H^\vee) = (F^{1-p}H)^\perp$.
\end{enumerate}
\end{theorem}

\begin{proof}
(1) By Theorem~\ref{mixed-hodge-structures:theorem-splitting}, $f_\CC$ maps $I^{p,q}(H)$ into $I^{p,q}(H')$. Let $y = f(x) \in W'_{k,\CC}$ and decompose $x = \sum x^{p,q}$ with $x^{p,q} \in I^{p,q}(H)$. Then
$y = \sum f(x^{p,q})$ is the decomposition of $y$, so $f(x^{p,q}) = 0$ whenever $p + q > k$, and $y = f(\sum_{p+q \leq k}x^{p,q}) \in f(W_{k,\CC})$. The same argument, with $p' \geq p$ in place of $p + q \leq k$,
proves strictness for $F$. Strictness for $W$ over $\QQ$ follows from strictness over $\CC$, because $f(W_k)$ and $W'_k \cap f(H)$ are $\QQ$-subspaces with the same complexification.

(2) Since $f_\CC$ preserves the bigradings, its kernel $K_\CC$, image $f(H_\CC)$ and cokernel $C_\CC$ are bigraded: $K_\CC = \bigoplus(K_\CC \cap I^{p,q}(H))$, $f(H_\CC) = \bigoplus f(I^{p,q}(H))$ and
$C_\CC = \bigoplus I^{p,q}(H')/f(I^{p,q}(H))$. Give $K$, $f(H)$ and $C$ the filtrations induced as subspaces or quotients. In each case $W_k$ and $F^p$ are the sums of the pieces with $p + q \leq k$,
respectively with first index $\geq p$. For $f(H)$ the filtrations induced from $H$ (quotient) and from $H'$ (subspace) agree, which is the strictness proved in (1). The bigradings satisfy the
conjugation condition of Theorem~\ref{mixed-hodge-structures:theorem-splitting}, because $K$, $f(H)$ and $C$ are defined over $\QQ$ and the components of a vector in a bigraded subspace lie in it. By the
converse part of Theorem~\ref{mixed-hodge-structures:theorem-splitting}, $K$, $f(H)$ and $C$ are mixed Hodge structures, and the canonical map $H/K \to f(H)$ is an isomorphism of mixed Hodge structures.

(3) If $0 \to H' \to H \to H'' \to 0$ is an exact sequence of mixed Hodge structures, then its complexification is exact and each map preserves the bigradings, so
$0 \to I^{p,q}(H') \to I^{p,q}(H) \to I^{p,q}(H'') \to 0$ is exact for each $(p, q)$. Since $\mathrm{Gr}^W_kH_\CC \cong \bigoplus_{p+q=k}I^{p,q}(H)$, the sequence of graded pieces is exact.

(4) Put $I^{p,q}(H \otimes H') = \bigoplus I^{a,b}(H) \otimes I^{a',b'}(H')$, the sum over $a + a' = p$, $b + b' = q$. This bigrading splits the filtrations $W$ and $F$ of the statement, and it satisfies the
conjugation condition because the error terms of the factors multiply to error terms. As $W_k(H \otimes H')$ is defined over $\QQ$, the converse part of Theorem~\ref{mixed-hodge-structures:theorem-splitting}
applies. For the dual, use $I^{p,q}(H^\vee) = (\bigoplus_{(p',q') \neq (-p,-q)}I^{p',q'}(H))^\perp$ in the same way.
\end{proof}

\begin{lemma}[Principle of weights]\label{mixed-hodge-structures:lemma-weights}
Let $f \colon H \to H'$ be a morphism of mixed Hodge structures.
\begin{enumerate}
\item If the weights of $H$ are all $< m$ and those of $H'$ are all $\geq m$, or the other way round, then $f = 0$.
\item If $H'$ is pure of weight $k$, then $f(H) = f(W_kH)$ and $f(W_{k-1}H) = 0$.
\end{enumerate}
\end{lemma}

\begin{proof}
(1) If the weights of $H$ are $< m$ and those of $H'$ are $\geq m$, then $H = W_{m-1}H$ and $f(H) \subseteq W'_{m-1}H' = 0$. In the other case $W_{m-1}H = 0$ and $H' = W'_{m-1}H'$, so strictness
(Theorem~\ref{mixed-hodge-structures:theorem-abelian}(1)) gives $f(H) = f(H) \cap W'_{m-1}H' = f(W_{m-1}H) = 0$.

(2) $W'_{k-1}H' = 0$ and $W'_kH' = H'$, so $f(W_{k-1}H) \subseteq W'_{k-1}H' = 0$, and strictness gives $f(W_kH) = W'_kH' \cap f(H) = f(H)$.
\end{proof}

\begin{remark}\label{mixed-hodge-structures:remark-principle}
Lemma~\ref{mixed-hodge-structures:lemma-weights} is used over and over below: a map compatible with mixed Hodge structures cannot move classes of one weight to classes of another. Its power comes from
the fact that the maps it is applied to (restrictions, Gysin maps, differentials of spectral sequences) are defined topologically, while the weights are defined algebro-geometrically. It has no
analogue for filtered vector spaces, where strictness fails (Counterexample~\ref{hodge-structures:counterexample-filtered}).
\end{remark}

\section{The logarithmic de Rham complex}\label{mixed-hodge-structures:section-log-de-rham}

\begin{definition}\label{mixed-hodge-structures:definition-log-forms}
Let $\bar X$ be a complex manifold of dimension $n$ and $D \subseteq \bar X$ a divisor with \emph{simple normal crossings}: $D = D_1 \cup \cdots \cup D_N$ with each $D_i$ a nonsingular hypersurface, and
near every point there are holomorphic coordinates $z_1, \ldots, z_n$ with $D = \{z_1 \cdots z_r = 0\}$. Put $X = \bar X \setminus D$ and let $j \colon X \to \bar X$ be the inclusion.
\begin{enumerate}
\item The sheaf $\Omega^p_{\bar X}(\log D)$ of \emph{logarithmic $p$-forms} is the $\mathcal{O}_{\bar X}$-submodule of $j_*\Omega^p_X$ generated locally by the wedge products of $p$ of the forms
$dz_1/z_1, \ldots, dz_r/z_r, dz_{r+1}, \ldots, dz_n$. It is locally free, and stable under $d$, which gives the \emph{logarithmic de Rham complex} $\Omega^\bullet_{\bar X}(\log D)$.
\item The \emph{weight filtration} $W_m\Omega^p_{\bar X}(\log D)$ is the subsheaf spanned by the forms $\alpha \wedge dz_{i_1}/z_{i_1} \wedge \cdots \wedge dz_{i_m}/z_{i_m}$ with $\alpha$ holomorphic: forms with at
most $m$ logarithmic factors.
\item $D^{(m)}$ is the disjoint union of the $m$-fold intersections $D_{i_1} \cap \cdots \cap D_{i_m}$, $i_1 < \cdots < i_m$, with $D^{(0)} = \bar X$. It is a complex manifold of dimension $n - m$, with the map
$a_m \colon D^{(m)} \to \bar X$.
\end{enumerate}
\end{definition}

\begin{lemma}\label{mixed-hodge-structures:lemma-residue}
The residue map $\alpha \wedge dz_{i_1}/z_{i_1} \wedge \cdots \wedge dz_{i_m}/z_{i_m} \mapsto \alpha|_{D_{i_1} \cap \cdots \cap D_{i_m}}$ induces an isomorphism of complexes
\[
\mathrm{Res}_m \colon \mathrm{Gr}^W_m\Omega^\bullet_{\bar X}(\log D) \longrightarrow (a_m)_*\Omega^{\bullet-m}_{D^{(m)}}.
\]
\end{lemma}

\begin{proof}[Sketch of proof]
It is easiest to construct the inverse. For $I = \{i_1 < \cdots < i_m\}$ write $D_I = D_{i_1} \cap \cdots \cap D_{i_m}$ and $dz_I/z_I = dz_{i_1}/z_{i_1} \wedge \cdots \wedge dz_{i_m}/z_{i_m}$. Given a
holomorphic form $\beta$ on $D_I$, extend it locally to a holomorphic form $\tilde\beta$ on $\bar X$ and send it to the class of $\tilde\beta \wedge dz_I/z_I$ in $\mathrm{Gr}^W_m$. The class does not depend on the
extension: two extensions differ by a combination of forms $z_i\gamma$ and $dz_i \wedge \gamma$ with $i \in I$, and $z_i\gamma \wedge dz_I/z_I$ and $dz_i \wedge \gamma \wedge dz_I/z_I$ have at most $m - 1$
logarithmic factors, since $z_i \cdot dz_i/z_i = dz_i$ and $dz_i \wedge dz_i/z_i = 0$. It does not depend on the coordinates: replacing $z_i$ by $uz_i$ with $u$ a unit changes $dz_i/z_i$ by the holomorphic
form $du/u$, which lowers the weight. The map is surjective by the definition of $W_m$, it commutes with $d$ because $d(dz_i/z_i) = 0$, and $\mathrm{Res}_m$ is its inverse. That it is injective, that is,
that a sum $\sum_I\tilde\beta_I \wedge dz_I/z_I$ lies in $W_{m-1}$ only if every $\beta_I$ vanishes, is a computation in the local basis $dz_I/z_I \wedge dz_J$, $I \subseteq \{1, \ldots, r\}$,
$J \subseteq \{r+1, \ldots, n\}$, which we omit \cite{deligne-hodge-2}.
\end{proof}

The residue depends on the ordering $i_1 < \cdots < i_m$ through a sign, and in the comparison with rational cohomology it comes with the factor $(2\pi i)^{-m}$, because $\frac{1}{2\pi i}\frac{dz}{z}$ is the
rational class of Example~\ref{mixed-hodge-structures:example-mhs}(2). This is why the Tate twist $(-m)$ appears below.

\begin{theorem}[Deligne]\label{mixed-hodge-structures:theorem-smooth}
Let $X$ be a nonsingular variety, and $\bar X$ a nonsingular compactification with $D = \bar X \setminus X$ a simple normal crossings divisor; such a compactification exists by Nagata's compactification
theorem and Hironaka's resolution of singularities \cite{hironaka-1964}.
\begin{enumerate}
\item The inclusion $\Omega^\bullet_{\bar X}(\log D) \to j_*\Omega^\bullet_X$ is a quasi-isomorphism, and $j_*\Omega^\bullet_X$ computes $Rj_*\CC_X$; hence
$\mathbb{H}^k(\bar X, \Omega^\bullet_{\bar X}(\log D)) \cong H^k(X; \CC)$.
\item Let $W_mH^k(X; \CC)$ be the image of $\mathbb{H}^k(\bar X, W_{m-k}\Omega^\bullet_{\bar X}(\log D))$ and $F^pH^k(X; \CC)$ the image of $\mathbb{H}^k(\bar X, \Omega^{\geq p}_{\bar X}(\log D))$. Then $W$ is
defined over $\QQ$, and $(H^k(X; \QQ), W, F)$ is a mixed Hodge structure, independent of the choice of $\bar X$ and functorial for morphisms of nonsingular varieties.
\item The weights of $H^k(X; \QQ)$ lie in $[k, 2k]$, and $W_kH^k(X; \QQ) = \operatorname{im}\big(H^k(\bar X; \QQ) \to H^k(X; \QQ)\big)$.
\item There is a spectral sequence, the \emph{weight spectral sequence},
\[
E_1^{-m,k+m} = H^{k-m}(D^{(m)}; \QQ)(-m) \Longrightarrow H^k(X; \QQ),
\]
whose terms are pure Hodge structures of weight $k + m$, whose differential $d_1$ is an alternating sum of Gysin maps, and which degenerates at $E_2$, with $E_2^{-m,k+m} = \mathrm{Gr}^W_{k+m}H^k(X; \QQ)$.
The Hodge-to-de Rham spectral sequence $E_1^{p,q} = H^q(\bar X, \Omega^p_{\bar X}(\log D)) \Rightarrow H^{p+q}(X; \CC)$ degenerates at $E_1$.
\end{enumerate}
\end{theorem}

This is \cite[Section 3]{deligne-hodge-2}; see also \cite[Chapter 4]{peters-steenbrink}. The weight spectral sequence is the spectral sequence of the filtered complex $(\Omega^\bullet_{\bar X}(\log D), W)$,
with $E_1$ computed by Lemma~\ref{mixed-hodge-structures:lemma-residue}. Once one knows that its differentials are morphisms of Hodge structures, degeneration at $E_2$ is the principle of weights: for
$r \geq 2$, $d_r \colon E_r^{-m,k+m} \to E_r^{-m+r,k+m-r+1}$ goes between subquotients of pure Hodge structures of weights $k + m$ and $k + m + 1 - r \neq k + m$, so it vanishes by
Lemma~\ref{mixed-hodge-structures:lemma-weights}(1).

\begin{example}\label{mixed-hodge-structures:example-punctured-curve}
Let $\bar C$ be a compact Riemann surface of genus $g$, $s \geq 1$, and $C = \bar C \setminus \{p_1, \ldots, p_s\}$. Then $D^{(1)} = \{p_1, \ldots, p_s\}$ and $D^{(2)} = \emptyset$, and the weight spectral
sequence reduces to the exact sequence of mixed Hodge structures
\[
0 \to H^1(\bar C; \QQ) \to H^1(C; \QQ) \xrightarrow{\ \mathrm{Res}\ } \QQ(-1)^{s} \xrightarrow{\ \Sigma\ } \QQ(-1) \to 0,
\]
where $\mathrm{Res}$ takes the residues at the punctures and $\Sigma$, the Gysin map $H^0(D^{(1)})(-1) \to H^2(\bar C)$, is the sum. So $W_1H^1(C) = H^1(\bar C)$ is pure of weight $1$ and
$\mathrm{Gr}^W_2H^1(C) \cong \QQ(-1)^{s-1}$; exactness at $\QQ(-1)^s$ is the residue theorem. For $\PP^1$ minus the three points $0, 1, \infty$, $H^1 \cong \QQ(-1)^2$ is spanned by the classes of
$\frac{1}{2\pi i}\frac{dz}{z}$ and $\frac{1}{2\pi i}\frac{dz}{z - 1}$.
\end{example}

\begin{counterexample}\label{mixed-hodge-structures:counterexample-odd-weight}
The weights of $H^k(X)$ for $X$ nonsingular need not include $k$ or $2k$, and cohomology of even degree can be pure of odd weight. Let $E \subseteq \PP^2$ be a smooth cubic curve and $X = \PP^2 \setminus E$,
an affine surface. The weight spectral sequence ($D^{(2)} = \emptyset$) gives the Gysin sequence
\[
H^0(E)(-1) \to H^2(\PP^2) \to H^2(X) \to H^1(E)(-1) \to H^3(\PP^2) = 0.
\]
The first map sends $1$ to $\mathrm{cl}(E) = 3h \neq 0$, so it is an isomorphism over $\QQ$, and $H^2(X; \QQ) \cong H^1(E; \QQ)(-1)$ is pure of weight $3$.
\end{counterexample}

\section{Singular varieties}\label{mixed-hodge-structures:section-singular}

\begin{theorem}[Deligne]\label{mixed-hodge-structures:theorem-singular}
For every variety $X$ and every $k$, $H^k(X; \QQ)$ and $H^k_c(X; \QQ)$ carry functorial mixed Hodge structures, extending those of Theorem~\ref{mixed-hodge-structures:theorem-smooth}. The long exact
sequence of a pair $(X, Z)$ with $Z$ closed, and the sequence $\cdots \to H^k_c(X \setminus Z) \to H^k_c(X) \to H^k_c(Z) \to \cdots$, are exact sequences of mixed Hodge structures, and cup products and
K\"unneth isomorphisms are morphisms of mixed Hodge structures. The weights of $H^k(X)$ lie in $[0, 2k]$; if $X$ is proper they lie in $[0, k]$, and if $X$ is nonsingular in $[k, 2k]$. If $X$ is proper and
$\pi \colon \tilde X \to X$ is a resolution of singularities, then $W_{k-1}H^k(X; \QQ) = \ker\big(\pi^* \colon H^k(X; \QQ) \to H^k(\tilde X; \QQ)\big)$.
\end{theorem}

This is \cite{deligne-hodge-3}, via simplicial resolutions of $X$ by nonsingular varieties; see also \cite[Chapter 5]{peters-steenbrink}. For proper $X$ with a resolution $\pi \colon \tilde X \to X$ that is an
isomorphism outside a closed subset $Z \subseteq X$ with preimage $E = \pi^{-1}(Z)$, there is an exact sequence of mixed Hodge structures
\[
\cdots \to H^{k-1}(E) \to H^k(X) \to H^k(\tilde X) \oplus H^k(Z) \to H^k(E) \to H^{k+1}(X) \to \cdots,
\]
the Mayer--Vietoris sequence of the proper modification, which computes the mixed Hodge structure of $X$ inductively on the dimension.

\begin{example}\label{mixed-hodge-structures:example-nodal}
Let $C$ be the nodal cubic, $\pi \colon \PP^1 \to C$ its normalisation, $Z$ the node and $E = \pi^{-1}(Z) = \{0, \infty\}$. The sequence above reads
\[
0 \to H^0(C) \to H^0(\PP^1) \oplus H^0(Z) \to H^0(E) \to H^1(C) \to H^1(\PP^1) \oplus H^1(Z) = 0.
\]
The map $\QQ^2 \to \QQ^2$, $(a, b) \mapsto (a - b, a - b)$, has a one-dimensional cokernel, so $H^1(C) \cong \QQ(0)$, a quotient of the pure structure $H^0(E)$ of weight $0$. Accordingly
$W_0H^1(C) = \ker(H^1(C) \to H^1(\PP^1)) = H^1(C)$, as Theorem~\ref{mixed-hodge-structures:theorem-singular} predicts, and $H^1(C)$ has no Hodge decomposition of weight $1$.
\end{example}

\section{Applications of weights}\label{mixed-hodge-structures:section-applications}

\begin{corollary}\label{mixed-hodge-structures:corollary-restriction}
Let $\bar X$ be a nonsingular projective variety of dimension $n$, $U \subseteq \bar X$ a nonempty open subset, and $Z = \bar X \setminus U$ with irreducible components $Z_1, \ldots, Z_N$ of codimensions
$c_1, \ldots, c_N$. Let $g_i \colon \tilde Z_i \to \bar X$ be resolutions of singularities of the $Z_i$ composed with the inclusions.
\begin{enumerate}
\item The image of the restriction $H^k(\bar X; \QQ) \to H^k(U; \QQ)$ is $W_kH^k(U; \QQ)$.
\item The kernel of the restriction is $\sum_ig_{i*}H^{k-2c_i}(\tilde Z_i; \QQ)(-c_i)$, the sum of the images of the Gysin maps.
\end{enumerate}
\end{corollary}

\begin{proof}[Sketch of proof]
(1) By resolution of singularities there is a proper birational morphism $\pi \colon \tilde X \to \bar X$ from a nonsingular projective variety that is an isomorphism over $U$ and such that
$\tilde X \setminus \pi^{-1}(U)$ is a simple normal crossings divisor. By Theorem~\ref{mixed-hodge-structures:theorem-smooth}(3), the image of $H^k(\tilde X) \to H^k(U)$ is $W_kH^k(U)$. The image of
$H^k(\bar X)$ is contained in it, because restriction from $\bar X$ factors through $\pi^*$. Conversely, for $\alpha \in H^k(\tilde X)$ the class $\pi_*\alpha \in H^k(\bar X)$, defined by Poincar\'e duality,
restricts on $U$ to $\alpha|_U$, since $\pi$ is an isomorphism over $U$. So the two images agree.

(2) The Gysin images restrict to zero on $U$, since the classes are supported on $Z$. Conversely, the kernel is the image of $H^k_Z(\bar X; \QQ) \cong H^{BM}_{2n-k}(Z; \QQ)(-n)$ (Poincar\'e--Alexander
duality), by a morphism of mixed Hodge structures into the pure structure $H^k(\bar X)$ of weight $k$. By Lemma~\ref{mixed-hodge-structures:lemma-weights}(2) the kernel is the image of the part of weight
$\leq k$, and Deligne shows that the lowest weight part of the Borel--Moore homology of $Z$ is the image of the homology of the resolutions $\tilde Z_i$ \cite{deligne-hodge-3}. This last
identification is the verification we omit.
\end{proof}

\begin{theorem}[Global invariant cycle theorem]\label{mixed-hodge-structures:theorem-global-invariant}
Let $f \colon \mathcal{X} \to S$ be a smooth projective morphism of nonsingular varieties with $S$ connected, $\bar{\mathcal{X}}$ a nonsingular projective compactification of $\mathcal{X}$, and $t \in S$ with
fibre $X_t$. Then
\[
\operatorname{im}\big(H^k(\bar{\mathcal{X}}; \QQ) \to H^k(X_t; \QQ)\big) = \operatorname{im}\big(H^k(\mathcal{X}; \QQ) \to H^k(X_t; \QQ)\big) = H^k(X_t; \QQ)^{\pi_1(S, t)},
\]
and this subspace is a sub-Hodge structure of $H^k(X_t; \QQ)$.
\end{theorem}

\begin{proof}
The Leray spectral sequence $E_2^{a,b} = H^a(S, R^bf_*\QQ) \Rightarrow H^{a+b}(\mathcal{X}; \QQ)$ of a smooth projective morphism degenerates at $E_2$ \cite{deligne-1968}. Its edge map
$H^k(\mathcal{X}) \to E_\infty^{0,k} \subseteq E_2^{0,k} = H^0(S, R^kf_*\QQ) \subseteq H^k(X_t)$ is the restriction to the fibre, and by degeneration $E_\infty^{0,k} = E_2^{0,k}$, so the restriction is
surjective onto $H^0(S, R^kf_*\QQ)$. Since $R^kf_*\QQ$ is a local system (Theorem~\ref{vhs:theorem-ehresmann}) and $S$ is connected, $H^0(S, R^kf_*\QQ)$ is the space of $\pi_1(S, t)$-invariants in the stalk
$H^k(X_t; \QQ)$. This proves the second equality.

The restriction $r \colon H^k(\mathcal{X}; \QQ) \to H^k(X_t; \QQ)$ is induced by the inclusion of nonsingular varieties $X_t \to \mathcal{X}$, so it is a morphism of mixed Hodge structures
(Theorem~\ref{mixed-hodge-structures:theorem-smooth}(2)), and its target is pure of weight $k$ because $X_t$ is nonsingular and projective. By Lemma~\ref{mixed-hodge-structures:lemma-weights}(2),
$r(H^k(\mathcal{X})) = r(W_kH^k(\mathcal{X}))$, and $W_kH^k(\mathcal{X})$ is the image of $H^k(\bar{\mathcal{X}})$ by Corollary~\ref{mixed-hodge-structures:corollary-restriction}(1). Since restriction from
$\bar{\mathcal{X}}$ to $X_t$ factors through $\mathcal{X}$, this proves the first equality. Finally, the image of $H^k(\bar{\mathcal{X}}) \to H^k(X_t)$ is the image of a morphism of pure Hodge structures of
weight $k$, hence a sub-Hodge structure (Theorem~\ref{hodge-structures:theorem-abelian}).
\end{proof}

\begin{remark}\label{mixed-hodge-structures:remark-fixed-part}
Theorem~\ref{mixed-hodge-structures:theorem-global-invariant} is the proof of Theorem~\ref{vhs:theorem-fixed-part} promised in Chapter~\ref{vhs}. The only input not proved in this book is the degeneration
of the Leray spectral sequence, which Deligne deduces from the hard Lefschetz theorem (Chapter~\ref{lefschetz-theorems}) applied to the fibres \cite{deligne-1968}. The coniveau filtration is handled in
the same way: $N^cH^k(X)$ is a sum of kernels of restrictions $H^k(X) \to H^k(X \setminus Z)$, which by Corollary~\ref{mixed-hodge-structures:corollary-restriction}(2) are images of Gysin maps and so
sub-Hodge structures (Proposition~\ref{hodge-conjecture:proposition-coniveau}).
\end{remark}

\begin{proposition}\label{mixed-hodge-structures:proposition-localisation}
Keep the notation of Corollary~\ref{mixed-hodge-structures:corollary-restriction}, and let $\alpha \in \mathrm{Hdg}^p(\bar X)$. Suppose that
\begin{enumerate}
\item there is an algebraic cycle $V$ on $\bar X$ with $\QQ$-coefficients such that $\alpha|_U = \mathrm{cl}(V)|_U$ in $H^{2p}(U; \QQ)$, and
\item for each $i$ with $c_i \leq p$, the Hodge conjecture holds for $\tilde Z_i$ in codimension $p - c_i$.
\end{enumerate}
Then $\alpha$ is algebraic.
\end{proposition}

\begin{proof}
Put $\beta = \alpha - \mathrm{cl}(V)$. It is a Hodge class, by Corollary~\ref{cycle-class:corollary-hodge-classes}, and $\beta|_U = 0$. By Corollary~\ref{mixed-hodge-structures:corollary-restriction}(2),
$\beta$ lies in the image of
\[
g_* = \sum_ig_{i*} \colon A = \bigoplus_iH^{2p-2c_i}(\tilde Z_i; \QQ)(-c_i) \longrightarrow H^{2p}(\bar X; \QQ),
\]
where the summands with $c_i > p$ are zero. Each Gysin map $g_{i*}$ is a morphism of Hodge structures of weight $2p$: it is Poincar\'e dual to pushforward in homology, and it maps forms of type $(a, b)$
to currents of type $(a + c_i, b + c_i)$. The source $A$ is a direct sum of polarisable Hodge structures of weight $2p$ (Theorem~\ref{hodge-structures:theorem-cohomology-polarised}). By
Theorem~\ref{hodge-structures:theorem-semisimple}, $\ker g_*$ has a complement $A'$ in $A$ that is a sub-Hodge structure, and $g_*|_{A'} \colon A' \to g_*(A)$ is an isomorphism of Hodge structures onto the
sub-Hodge structure $g_*(A)$. The preimage $\gamma = (\gamma_i) \in A'$ of $\beta$ is then a Hodge class, since isomorphisms of Hodge structures preserve Hodge classes
(Proposition~\ref{hodge-structures:proposition-hodge-classes-morphisms}); that is, $\gamma_i \in \mathrm{Hdg}^{p-c_i}(\tilde Z_i)$. By hypothesis (2), $\gamma_i = \mathrm{cl}(V_i)$ for cycles $V_i$ on $\tilde Z_i$
with $\QQ$-coefficients. By compatibility of the cycle class with proper pushforward (Theorem~\ref{cycle-class:theorem-cycle-class}(2)), $\beta = \sum_ig_{i*}\mathrm{cl}(V_i) = \mathrm{cl}(\sum_ig_{i*}V_i)$.
Hence $\alpha = \mathrm{cl}(V + \sum_ig_{i*}V_i)$ is algebraic.
\end{proof}

\begin{corollary}\label{mixed-hodge-structures:corollary-localisation}
Assume that the Hodge conjecture holds for all nonsingular projective varieties of dimension $< n$, and let $\bar X$ be an irreducible nonsingular projective variety of dimension $n$ and
$U \subseteq \bar X$ a nonempty open subset. Then a Hodge class $\alpha \in \mathrm{Hdg}^p(\bar X)$ is algebraic if and only if $\alpha|_U$ is the restriction of the class of an algebraic cycle on $\bar X$.
\end{corollary}

\begin{proof}
One direction is trivial. For the other, since $\bar X$ is irreducible and $U$ is nonempty, every component of $\bar X \setminus U$ has codimension $c_i \geq 1$, so $\dim\tilde Z_i < n$ and hypothesis (2) of
Proposition~\ref{mixed-hodge-structures:proposition-localisation} holds.
\end{proof}

\begin{corollary}\label{mixed-hodge-structures:corollary-birational}
Assume that the Hodge conjecture holds for all nonsingular projective varieties of dimension $< n$, and let $\bar X$ and $\bar X'$ be birational irreducible nonsingular projective varieties of dimension $n$.
Then the Hodge conjecture holds for $\bar X$ if and only if it holds for $\bar X'$.
\end{corollary}

\begin{proof}
Choose nonempty open subsets $U \subseteq \bar X$ and $U' \subseteq \bar X'$ and an isomorphism $\phi \colon U' \to U$; we identify $U$ and $U'$ through $\phi$. Suppose the Hodge conjecture holds for
$\bar X'$, and let $\alpha \in \mathrm{Hdg}^p(\bar X)$. The restriction $\alpha|_U$ lies in $W_{2p}H^{2p}(U; \QQ)$ by Corollary~\ref{mixed-hodge-structures:corollary-restriction}(1), which is a pure Hodge
structure of weight $2p$, since the weights of $H^{2p}(U)$ are $\geq 2p$ (Theorem~\ref{mixed-hodge-structures:theorem-smooth}(3)). As the restriction is a morphism of mixed Hodge structures, $\alpha|_U$ is
rational and lies in $F^p$, so it is a Hodge class of this pure structure (Proposition~\ref{hodge-structures:proposition-filtration}: in weight $2p$, $F^p \cap \overline{F^p}$ is the $(p,p)$-part). The restriction
$r' \colon H^{2p}(\bar X'; \QQ) \to W_{2p}H^{2p}(U; \QQ)$ is a surjective morphism of pure Hodge structures of weight $2p$, by Corollary~\ref{mixed-hodge-structures:corollary-restriction}(1) applied to $\bar X'$.
As in the proof of Proposition~\ref{mixed-hodge-structures:proposition-localisation}, semisimplicity (Theorem~\ref{hodge-structures:theorem-semisimple}) gives a Hodge class $\alpha' \in \mathrm{Hdg}^p(\bar X')$ with
$r'(\alpha') = \alpha|_U$. By hypothesis $\alpha' = \mathrm{cl}(V')$ for a cycle $V'$ on $\bar X'$ with $\QQ$-coefficients. Let $V$ be the closure in $\bar X$ of the cycle $V'|_U$. Since the cycle class
commutes with restriction to open subsets (Theorem~\ref{cycle-class:theorem-cycle-class}(2)), $\mathrm{cl}(V)|_U = \mathrm{cl}(V'|_U) = \alpha'|_U = \alpha|_U$. By
Corollary~\ref{mixed-hodge-structures:corollary-localisation}, $\alpha$ is algebraic. The converse follows by symmetry.
\end{proof}

\begin{remark}\label{mixed-hodge-structures:remark-integral}
The arguments of Proposition~\ref{mixed-hodge-structures:proposition-localisation} and Corollary~\ref{mixed-hodge-structures:corollary-birational} use semisimplicity to lift Hodge classes along surjections,
and this is a statement about $\QQ$-structures. Over $\ZZ$ an integral Hodge class in the image of a morphism need not be the image of an integral Hodge class, and the integral analogues must be proved
by other means: for classes of degree $4$, the birational invariance of the defect $Z^4(X)$ is a theorem of Colliot-Th\'el\`ene and Voisin (Theorem~\ref{failures:theorem-what-survives}).
\end{remark}

\section{Exercises}\label{mixed-hodge-structures:section-exercises}

\begin{exercise}\label{mixed-hodge-structures:exercise-cstar}
Compute the mixed Hodge structure on $H^\bullet((\CC^*)^n; \QQ)$.
\end{exercise}

\begin{solution}
By the K\"unneth formula, which is compatible with mixed Hodge structures (Theorem~\ref{mixed-hodge-structures:theorem-singular}), and Example~\ref{mixed-hodge-structures:example-mhs}(2),
$H^k((\CC^*)^n) \cong \Lambda^k(\QQ(-1)^n) \cong \QQ(-k)^{\binom{n}{k}}$, pure of weight $2k$ and type $(k,k)$, spanned by the products of the classes $\frac{1}{2\pi i}\frac{dz_j}{z_j}$. The weights are the largest
allowed by Theorem~\ref{mixed-hodge-structures:theorem-smooth}(3).
\end{solution}

\begin{exercise}\label{mixed-hodge-structures:exercise-kummer}
Show that the Kummer structure $K_q$ of Example~\ref{mixed-hodge-structures:example-mhs}(4) is a mixed Hodge structure with $\mathrm{Gr}^W_{-2}K_q \cong \QQ(1)$ and $\mathrm{Gr}^W_0K_q \cong \QQ(0)$, and that
$K_q \cong \QQ(1) \oplus \QQ(0)$ if and only if $q$ is a root of unity.
\end{exercise}

\begin{solution}
The only nonzero graded pieces are $\mathrm{Gr}^W_{-2} = \QQ e_1$ and $\mathrm{Gr}^W_0 = H/\QQ e_1$. On $\mathrm{Gr}^W_{-2}$ the induced filtration is $F^{-1} = \CC e_1$ and $F^0 = F^0 \cap \CC e_1 = 0$, since
$e_0 - \lambda e_1 \notin \CC e_1$; this is the Hodge filtration of type $(-1,-1)$, that of $\QQ(1)$ (identify $e_1$ with $2\pi i$). On $\mathrm{Gr}^W_0$ the image of $F^0$ is everything and $F^1 = 0$, which is
type $(0,0)$. So $K_q$ is a mixed Hodge structure. An isomorphism $K_q \cong \QQ(1) \oplus \QQ(0)$ maps $W_{-2}K_q$ onto $\QQ(1)$, so the preimage of $\QQ(0)$ is a complement of $W_{-2}$ that is a
sub-mixed Hodge structure; by Counterexample~\ref{mixed-hodge-structures:counterexample-no-splitting} one exists if and only if $\lambda = \log q/2\pi i \in \QQ$. Writing $\lambda = a/b$ with $b \geq 1$, this
says $q^b = e^{2\pi ia} = 1$; conversely if $q^b = 1$ then $\log q \in 2\pi i\frac{1}{b}\ZZ$.
\end{solution}

\begin{exercise}\label{mixed-hodge-structures:exercise-hodge-number-bounds}
Let $X$ be a nonsingular variety. Show that $h^{p,q}(H^k(X)) = 0$ unless $0 \leq p, q \leq k$ and $k \leq p + q \leq 2k$.
\end{exercise}

\begin{solution}
The condition on $p + q$ is Theorem~\ref{mixed-hodge-structures:theorem-smooth}(3). The complex $\Omega^{\geq k+1}_{\bar X}(\log D)$ is zero in degrees $\leq k$, so its hypercohomology vanishes in degree $k$, and
$F^{k+1}H^k(X; \CC) = 0$. Hence $F^{k+1}$ vanishes on every $\mathrm{Gr}^W_m$ and $h^{p,q} = 0$ for $p > k$; by the symmetry $h^{p,q} = h^{q,p}$ of pure Hodge structures, also for $q > k$. Finally
$F^0 = H_\CC$ because $\Omega^{\geq 0}(\log D)$ is the whole complex, so $h^{p,q} = 0$ for $p < 0$, and for $q < 0$ by symmetry.
\end{solution}

\begin{exercise}\label{mixed-hodge-structures:exercise-elliptic-punctured}
Let $E$ be an elliptic curve with origin $0$ and $a \in E$, $a \neq 0$. Show that $H^1(E \setminus \{0\})$ is pure of weight $1$, that $H^1(E \setminus \{0, a\})$ has $\mathrm{Gr}^W_1 \cong H^1(E)$ and $\mathrm{Gr}^W_2 \cong \QQ(-1)$, and compute the
Hodge numbers of both.
\end{exercise}

\begin{solution}
By Example~\ref{mixed-hodge-structures:example-punctured-curve} with $g = 1$: for $s = 1$, $\mathrm{Gr}^W_2 \cong \QQ(-1)^0 = 0$, so $H^1(E \setminus \{0\}) = H^1(E)$, with $h^{1,0} = h^{0,1} = 1$. For $s = 2$,
$\mathrm{Gr}^W_1 = H^1(E)$ and $\mathrm{Gr}^W_2 \cong \QQ(-1)$, so $h^{1,0} = h^{0,1} = h^{1,1} = 1$ and $\dim H^1 = 3$. (Whether this extension splits depends on $a$: it splits if and only if $a$ is a torsion
point of $E$, by Carlson's description of extensions \cite{carlson-1980}; we do not prove this.)
\end{solution}
